% ---------------------------------------------------------------------------
%  QML Lab research project 11: Quantum Reservoir Computing for Time Series
%  Build:  latexmk -pdf main.tex   (needs qmlproject.sty, harvardqml.bst, references.bib)
% ---------------------------------------------------------------------------
\documentclass[11pt,a4paper]{article}
\usepackage{qmlproject}
\qmlheader{Project 11: Quantum reservoir computing}

\hypersetup{pdftitle={QML Lab project 11: Quantum Reservoir Computing for Time Series}}

\begin{document}

\begin{center}
  {\small\textcolor{black!55}{QML Lab \textbullet\ research project 11}}\\[0.5em]
  {\LARGE\bfseries Quantum Reservoir Computing for Time Series}\\[0.6em]
  {\large A fixed random quantum system as a reservoir, with only a linear read-out trained.}\\[0.6em]
  {\small Living document \textbullet\ last updated \today}
\end{center}

\begin{projectcard}
\small
\begin{tabularx}{\linewidth}{@{}>{\bfseries}l X >{\bfseries}l X@{}}
Quadrant  & \CQ\ (classical data, quantum processor) & Track & Hybrid \\
Difficulty & Medium & Tools & PennyLane, scikit-learn \\
Builds on & Modules 2 and 4 of the QML Lab & Workspace & \path{qml_projects/quantum-reservoir-time-series/} \\
\end{tabularx}
\end{projectcard}

\paragraph{Summary.}
Reservoir computing avoids training the hard part of a recurrent model: a fixed dynamical system
mixes the input history, and only a linear read-out is fitted. This project uses a small, fixed
quantum circuit as the reservoir, feeds a time series into it step by step, and forecasts the series
from measured expectation values. It compares the quantum reservoir with a classical echo-state network
of the same size on standard forecasting and memory tasks.

% ===========================================================================
\section{Background and key concepts}
% ===========================================================================
\subsection{Reservoir computing}
An echo-state network \citep{jaeger2004} drives a fixed random recurrent network with the input
$u_t$ and fits only the output weights. With features $x_t\in\R^m$ collected in a matrix $X$ and
targets $Y$, the read-out is ridge regression,
\begin{equation}
  W=YX^\top\big(XX^\top+\lambda I\big)^{-1},\qquad \hat y_t=Wx_t ,
\end{equation}
which trains in closed form. The reservoir needs \emph{fading memory}: recent inputs must matter more
than old ones.

\subsection{Quantum reservoirs}
\citet{fujii2017} proposed quantum reservoirs: the input is injected into one qubit at every step,
the system evolves under a fixed Hamiltonian, and expectation values form the features. In a circuit
version,
\begin{equation}
  \rho_{t+1}=\mathcal{U}\Big(\ket{u_t}\!\bra{u_t}\otimes\Tr_1[\rho_t]\Big)\,\mathcal{U}^\dagger,\qquad
  x_t=\big(\langle Z_i\rangle_t,\ \langle Z_iZ_j\rangle_t\big)_{i<j},
\end{equation}
where $\ket{u_t}=R_Y(u_t)\ket{0}$ replaces qubit 1 at each step (the partial reset gives fading
memory). The exponentially large state space promises rich dynamics from few qubits; the review of
\citet{mujal2021} discusses the opportunities and the cost of measurement.

\subsection{Benchmarks}
The Mackey--Glass system \citep{mackey1977},
$\dot x(t)=\beta\,x(t-\tau)/\big(1+x(t-\tau)^{10}\big)-\gamma x(t)$ with $\beta=0.2$, $\gamma=0.1$ and
$\tau=17$, is chaotic and a standard forecasting test. Memory capacity
$\mathrm{MC}=\sum_{d\ge1}r^2\big(u_{t-d},\hat y^{(d)}_t\big)$ measures how many past inputs the
reservoir can reconstruct.

% ===========================================================================
\section{Research questions}
% ===========================================================================
The project is open-ended: these questions are the starting point, not the finish line. Add new
ones to the README and the journal as they come up.

\begin{rqlist}
  \item How does forecast error depend on reservoir size and on which observables are measured?
  \item Does memory of past inputs fade too fast or too slowly, and how can the reset strength tune it?
  \item How does the quantum reservoir compare with a classical echo-state network with the same number of features?
\end{rqlist}

% ===========================================================================
\section{Method}
% ===========================================================================
\subsection{Data}
A noisy sine wave (sanity check), the Mackey--Glass series sampled at unit time steps, and
i.i.d.\ uniform inputs for memory capacity. 100 washout steps, 1000 training steps, 500 test steps.

\subsection{Models}
$n=4$--$8$ qubits simulated as density matrices (\texttt{default.mixed}); fixed random rotations and
CNOTs as $\mathcal U$; input injection by reset and $R_Y(u_t)$ on qubit 1; single- and two-qubit $Z$
expectation values as features (exact, then from $M$ shots); ridge read-out.

\subsection{Baselines}
Echo-state network with the same number of read-out features; linear autoregressive model; a small
LSTM as a strong upper reference.

\subsection{Metrics}
\begin{itemize}
  \item Normalised root mean squared error (NRMSE) on test data.
  \item Memory capacity curve $r^2$ versus delay $d$.
  \item Robustness to the number of shots $M$.
\end{itemize}

% ===========================================================================
\section{Planned experiments}
% ===========================================================================
Each experiment becomes one \texttt{[[experiments]]} entry in the workspace, so its runs, parameters
and metrics are stored and compared in the portal.

\begin{table}[h]
\centering\small
\begin{tabularx}{\linewidth}{@{}l>{\raggedright}X>{\raggedright}X>{\raggedright\arraybackslash}X@{}}
\toprule
ID & Experiment & What varies & Main metrics \\
\midrule
  E1 & Sine forecast & reservoir size $n$ & NRMSE \\
  E2 & Mackey--Glass & quantum vs echo-state network & NRMSE one- and multi-step \\
  E3 & Features & single vs two-qubit observables & NRMSE \\
  E4 & Memory & delay $d=1..20$ & memory capacity \\
  E5 & Shot noise & $M\in\{100,1000,10^4\}$ & NRMSE \\
\bottomrule
\end{tabularx}
\caption{Experiment plan. Results go into the journal and the results log (\cref{sec:log}).}
\end{table}

% ===========================================================================
\section{Running it in the QML Lab}
% ===========================================================================
The workspace lives in \path{qml_projects/quantum-reservoir-time-series/}. Scripts import three helpers from
\path{qml_lab/plots.py}:
\begin{itemize}
  \item \texttt{params()} reads the parameters from the portal's run form (defaults when run by hand);
  \item \texttt{save()} stores figures with the run;
  \item \texttt{log\_metrics()} records numbers that appear in the run table.
\end{itemize}

\begin{lstlisting}
qml_projects/quantum-reservoir-time-series/
  project.toml          metadata + experiments
  README.md             goal, open questions, ideas
  experiments/          runnable scripts
  docs/main.tex         this document
\end{lstlisting}

First experiment: This experiment is \emph{planned}: write the script, add the block below to \texttt{project.toml}, then run \texttt{python manage.py sync\_projects}.

\begin{lstlisting}[language={}]
[[experiments]]
slug = "mackey-glass"
title = "Quantum reservoir forecast of Mackey--Glass"
script = "experiments/e01_mackey_glass.py"
figure = "forecast.png"

[experiments.params]
n_qubits = 6
washout = 100
train_steps = 1000
ridge = 1e-06
shots = 0
task = "mackey_glass"
seed = 0
\end{lstlisting}

\begin{lstlisting}[language=bash]
# from the portal folder: pick up changes to project.toml
python manage.py sync_projects
cd qml_projects/quantum-reservoir-time-series
# run by hand with the default parameters
PYTHONPATH=../../qml_lab python experiments/e01_mackey_glass.py
\end{lstlisting}

% ===========================================================================
\section{Quality bar}
% ===========================================================================
\begin{itemize}
  \item A classical baseline with a similar number of trainable parameters, tuned with the same care \citep{bowles2024}.
  \item At least five random seeds per configuration, reported as mean $\pm$ standard deviation.
  \item How results change with the number of qubits.
  \item At least one run under a realistic noise model.
  \item Give the echo-state network exactly as many read-out features as the quantum reservoir.
\end{itemize}

% ===========================================================================
\section{Open directions}
% ===========================================================================
\begin{itemize}
  \item Closed-loop (generative) forecasting with feedback.
  \item Quantum extreme learning machines, which drop the memory and process one input at a time.
  \item Hardware with mid-circuit measurement and reset.
\end{itemize}

% ===========================================================================
\section{Risks and limitations}
% ===========================================================================
\begin{itemize}
  \item On hardware, measuring destroys the reservoir state; each time step may require re-running the input history.
  \item For $n\le8$ qubits the reservoir is easy to simulate classically; any advantage must come from scaling.
  \item Results depend strongly on reservoir and ridge hyperparameters.
\end{itemize}

% ===========================================================================
\section{Results log}
\label{sec:log}
% ===========================================================================
Summarise findings here as they become stable; day-to-day notes belong in the portal journal.

\begin{table}[h]
\centering\small
\begin{tabularx}{\linewidth}{@{}l l >{\raggedright}X >{\raggedright\arraybackslash}X@{}}
\toprule
Date & Run & Finding & Next step \\
\midrule
 & & & \\
 & & & \\
 & & & \\
\bottomrule
\end{tabularx}
\end{table}

\bibliographystyle{harvardqml}
\bibliography{references}

\end{document}
